\documentclass[aspectratio=169]{beamer}
\usetheme{default}
\usecolortheme{default}
\title{Beamer Presentation}
\author{Generated}

\title[short title]{Long title}
\author[short author family name split by ',']{
    author1 \inst{1,2} \and
    author2 \inst{1}
}
\institute[short notation of institute; split by ',']{
    \inst{1} inst1 %
    \inst{2} inst2  %
}

\date{\today}

\usepackage[utf8]{inputenc}
\usepackage{amsmath}
\usepackage{graphicx}
\usepackage{hyperref}


\begin{document}


% create the title frame here


% BEGIN_DEMO_FRAME
% END_DEMO_FRAME
% BEGIN_FRAMES
\begin{frame}{Deep Learning for Natural Language Processing: Recent Advances and Applications}
\centering
\vspace{1cm}
\Huge\textbf{Deep Learning for Natural Language Processing: Recent Advances and Applications}

\vspace{1cm}
\normalsize
\begin{center}
\textbf{Authors:} John Doe, Jane Smith, Alex Lee \\
\textbf{Affiliation:} Department of Computer Science, University of Example \\
\small{Contact: john.doe@university.edu, jane.smith@university.edu}
\end{center}

\vspace{1cm}
\small
\textit{This work was supported by the National Science Foundation under Grant No. 123456.}
\end{frame}

\begin{frame}{Introduction}
Natural Language Processing (NLP) has witnessed remarkable progress in recent years, primarily driven by advances in deep learning. From machine translation to question answering, deep learning models have achieved unprecedented performance across a wide range of NLP tasks.

The contributions of this paper are threefold:
\begin{enumerate}
\item We provide a comprehensive survey of modern deep learning architectures for NLP
\item We present empirical comparisons of key models on standard benchmarks
\item We discuss practical deployment considerations and future research directions
\end{enumerate}
\end{frame}

\begin{frame}{Background: Neural Networks for NLP}
\begin{itemize}
  \item Recurrent Neural Networks (RNNs) and Long Short-Term Memory (LSTM) networks became popular for sequence modeling tasks.
  \item However, these architectures faced challenges with long-range dependencies.
\end{itemize}
\end{frame}

\begin{frame}{Transformer Architecture: Self-Attention}
\begin{block}{Self-Attention Mechanism}
Self-attention allows the model to weigh the importance of different tokens when processing each position. The attention function can be described as:
\end{block}

\begin{equation}
\text{Attention}(Q, K, V) = \text{softmax}\left(\frac{QK^T}{\sqrt{d_k}}\right)V
\end{equation}

\begin{itemize}
  \item Enables dynamic weighting of token importance
  \item Computes context-aware representations
  \item Core to Transformer's effectiveness
\end{itemize}
\end{frame}

\begin{frame}{Transformer Architecture: Multi-Head Attention}
\begin{itemize}
  \item Multi-head attention extends the attention mechanism by employing multiple attention heads, allowing the model to learn different types of relationships.
\end{itemize}
\end{frame}

\begin{frame}{Large Language Models: Scaling Laws}
\begin{itemize}
  \item LLMs follow predictable scaling laws, where performance improves with model size, data, and compute.
\end{itemize}
\end{frame}

\begin{frame}{Large Language Models: Instruction Tuning}
\begin{itemize}
  \item Instruction tuning enables LLMs to follow natural language instructions, making them more versatile and user-friendly.
\end{itemize}
\end{frame}

\begin{frame}{Experimental Results: GLUE Benchmark}
\begin{itemize}
    \item We evaluate several models on standard benchmarks. Table~\ref{tab:results} summarizes our findings.
\end{itemize}

\begin{table}[htbp]
\centering
\caption{Performance on GLUE benchmark}
\small
\setlength{\tabcolsep}{3pt}
\renewcommand{\arraystretch}{0.8}
\begin{tabular}{|l|c|c|}
\hline
Model & Accuracy & F1 Score \\
\hline
BERT-base & 89.3 & 88.7 \\
BERT-large & 91.5 & 91.2 \\
GPT-3 & 93.8 & 93.5 \\
GPT-4 & 95.2 & 95.0 \\
\hline
\end{tabular}
\label{tab:results}
\end{table}
\end{frame}

\begin{frame}{Discussion: Challenges}
\begin{itemize}
  \item Computational cost and environmental impact
  \item Bias and fairness considerations
  \item Interpretability and transparency
  \item Robustness to adversarial attacks
\end{itemize}
\end{frame}

\begin{frame}{Conclusion}
\begin{center}
\large
\textbf{Key Takeaways}
\end{center}

\vspace{1cm} % Increased spacing between title and content

\begin{itemize}
    \item We have surveyed recent advances in deep learning for NLP, focusing on transformer architectures and large language models.
    \item The field continues to evolve rapidly, with exciting opportunities for future research.
\end{itemize}

\vspace{1cm} % Additional spacing for visual hierarchy

\small
\textit{Sources: Vaswani et al. (2017), Mikolov et al. (2013)}
\end{frame}
% END_FRAMES
\end{document}